\scanpage{8}
Dear Murre,\\

Thank\unsure{s|x} you very much for your notes ont the tame fundamental group, which I at last finished reading. I see you wrote them with much care, and I am all the more sorry that \unsure{b|h}y my own fault, there is a number of misstatements which, I am affraid, will force you to do a serious recasting of the whole exposition. My notes definitely where too sketchy, and my oral explanations, I am affraid, partly wrong, which induced you into error a few times. Here the most serious drawbacks.

1.16. is false already when $H$ is the unit group and when there is a single $a$, say $a=b^n$, $Y'=Y\ T/(T^n-a)$. Then $Y=S$, and a morphism\deleted{\unsure{s|x}} of $Y$ into $Y'$ compatible with $H=e\quad H'$ is just a section\raisebox{-.45ex}{\deleted{\unsure{r|x}}} of $Y'$, which exists indeed; however $H\quad H'$ is not surjective. \unsure{3|\handwritten{3}}.6. is equally false, as you see by the previous example, using the given section to define an $H'$-morphism $H'\quad Y'$ which is not an isomorphism\rlap{\handwritten{\tikz[overlay,baseline]{\draw[line width=.5pt](-.03em,-.12ex)--(.36em,1.7ex);}}}. As a consequence, the proof in your notes of 3.7. breaks down (as it uses 1.16) and so does the proof of 3.8. (I did not try to check 3.8. by some different proof).
% The spaces in Y T, H=e H', H H', and H' Y' are blank in the scan; no bracket or arrow is supplied.
% The diagonal handwritten stem through the point after isomorphism is reproduced inline. Its combined reading is an exclamation mark.
% The repeated numeral 3 is typewriter-first, with handwritten-state alternative at its first occurrence; its numerical value is not in doubt.

I am affraid \deleted{the proof of} 6.4. is false as stated, and that the statement is correct only if the $D_i$ are regular. Indeed, the end of then \deleted{proo} proof seemed to me very dubious; be careful that the intertia groups are determined only up to interior automorphism ! There is however a \inserted{(tautological)} generalization of the theorem for regular $D_i$, corresponding to the data of a single divisor $D$ with normal crossings, and a variant of the notion of tame ramification for such a divisor, by demanding that the coverings should be tamely ramified locally fot the étale topology for the family of local irreductible components of $D$; it is this notion of tameness which would seem more adapted to the situation of par.9 .
% The cancelled proo reaches the right edge; no terminal f is visible in this source image.
% archive sheet number: 7
